\documentclass[article]{jss}

% Load required packages
\usepackage{amsmath}
\usepackage{amssymb}
\usepackage{amsfonts}
\usepackage{graphicx}
\usepackage{url}
\usepackage{hyperref}
\usepackage{booktabs}
\usepackage{color}
\usepackage{listings}

% JSS article information
\author{
  Peter Cotton\\
  Microprediction\\
  \email{peter.cotton@gmail.com}\\
  \url{https://www.microprediction.org}
}

\title{\pkg{HumpDay}: Multi-Dimensional Thurstone Calibration for Context-Specific Optimizer Selection}

% Shortened title for header
\Plaintitle{HumpDay: Multi-Dimensional Thurstone Calibration}
\Shorttitle{\pkg{HumpDay}: Thurstone Optimizer Calibration}

% Abstract
\Abstract{
Selecting appropriate optimization algorithms remains a critical challenge, with existing approaches providing either global rankings that ignore problem context or requiring complex installations that limit accessibility. We present \pkg{HumpDay}, a browser-based platform that applies multi-dimensional Thurstone calibration to provide context-specific optimizer recommendations. Rather than single global rankings, our approach maps optimizer performance across three key dimensions: landscape characteristics (smooth, multimodal, rugged), problem dimensionality (low, medium, high), and computational budget (low, high). Using 432 benchmark comparisons across derivative-free optimizers, we demonstrate that no single optimizer dominates globally, but clear specialization patterns emerge. For example, L-BFGS-B excels on smooth high-dimensional problems with limited budgets (performance 0.806), while Powell's method dominates rugged landscapes with ample computational resources (0.917). The platform runs entirely in-browser via Pyodide, enabling zero-installation access and interactive exploration of these performance relationships. This multi-dimensional calibration approach transforms optimizer selection from guesswork to principled decision-making based on problem characteristics.
}

% Keywords
\Keywords{optimization, multi-dimensional ranking, Thurstone model, context-specific recommendations, browser computing, derivative-free optimization}

% Plainkeywords for PDF metadata
\Plainkeywords{optimization, multi-dimensional ranking, Thurstone model, context-specific recommendations, browser computing, derivative-free optimization}

% Package name
\Address{
  Peter Cotton\\
  Microprediction\\
  \email{peter.cotton@gmail.com}\\
  URL: \url{https://www.microprediction.org}
}

\begin{document}

\section{Introduction} \label{sec:intro}

The fundamental question facing practitioners is not ``Which optimizer is best?'' but rather ``Which optimizer should I use for \emph{my specific problem}?'' Traditional benchmarking approaches fail to address this practical need, providing either global rankings that ignore problem context or static comparisons on narrow test suites.

Consider a practitioner facing a 10-dimensional multimodal optimization problem with a budget of 100 function evaluations. Existing resources might indicate that Differential Evolution ``generally outperforms'' Nelder-Mead, but provide no guidance on whether this holds for their specific combination of dimensionality, landscape type, and computational constraints.

\subsection{The Context-Dependency Problem}

Recent analysis reveals that optimizer performance exhibits strong context dependency:
\begin{itemize}
\item L-BFGS-B achieves 0.806 relative performance on smooth high-dimensional problems but only 0.500 on rugged landscapes
\item Powell's method excels on rugged problems with high budgets (0.917) but underperforms on smooth low-budget scenarios
\item No single optimizer ranks first across all problem contexts
\end{itemize}

This context-dependency undermines traditional global ranking systems borrowed from chess (Elo ratings) or sports competitions, which assume consistent relative performance across all ``matches.''

\subsection{Limitations of Existing Approaches}

Current optimization benchmarking suffers from several critical limitations:

\textbf{Global Rankings:} Platforms like COCO/BBOB \citep{hansen2021coco} and IOH Profiler \citep{doerr2018ioh} provide valuable insights but aggregate performance across diverse problems into single rankings, obscuring context-specific patterns.

\textbf{Installation Barriers:} Complex software dependencies limit accessibility, particularly for educational use and rapid prototyping scenarios.

\textbf{Static Test Suites:} Fixed benchmark functions may not represent the distribution of real-world problems encountered by practitioners.

\textbf{Inappropriate Statistical Models:} Elo rating systems assume transitive performance relationships that do not hold in optimization contexts.

\subsection{Multi-Dimensional Thurstone Calibration}

This paper introduces multi-dimensional Thurstone calibration as a principled approach to context-specific optimizer recommendation. Instead of asking ``Which optimizer is best overall?'' we address three practical questions:

\begin{enumerate}
\item \textbf{Landscape Specialization:} How do optimizers perform across smooth, multimodal, and rugged function classes?
\item \textbf{Dimensionality Scaling:} Which optimizers maintain effectiveness as problem dimension increases?
\item \textbf{Budget Efficiency:} How does performance vary with computational resource availability?
\end{enumerate}

By mapping optimizer abilities across these three dimensions simultaneously, we create a practical decision framework that transforms algorithm selection from guesswork to principled choice.

\section{Multi-Dimensional Thurstone Methodology} \label{sec:methods}

\subsection{Theoretical Foundation}

Traditional optimizer comparison relies on pairwise performance metrics aggregated using Elo rating systems. However, optimization performance exhibits characteristics poorly modeled by Elo assumptions:

\textbf{Context Dependency:} Algorithm A may outperform B on smooth functions while B outperforms A on multimodal landscapes.

\textbf{Multi-Dimensional Performance:} Optimizer capabilities span multiple independent dimensions (landscape type, dimensionality, budget efficiency).

\textbf{Continuous Outcomes:} Unlike binary chess outcomes, optimization produces continuous performance measures.

We employ the Thurstone model \citep{thurstone1927law, bradley1952rank} extended to multiple context dimensions:

\begin{equation}
P(\text{Algorithm } i \text{ beats } j | \text{context } \mathbf{c}) = \Phi\left(\frac{\theta_{i}(\mathbf{c}) - \theta_{j}(\mathbf{c})}{\sigma}\right)
\end{equation}

where $\theta_{i}(\mathbf{c})$ represents algorithm $i$'s latent ability in context $\mathbf{c} = (\text{landscape}, \text{dimension}, \text{budget})$, and $\Phi$ is the standard normal CDF.

\subsection{Context Dimension Definitions}

\textbf{Landscape Type Dimension:}
\begin{itemize}
\item \emph{Smooth:} Functions with few local optima (Rosenbrock, Zakharov)
\item \emph{Multimodal:} Functions with many local optima (Rastrigin, Griewank)
\item \emph{Rugged:} Functions with irregular, noisy landscapes (Shekel, combinations)
\end{itemize}

\textbf{Dimensionality Dimension:}
\begin{itemize}
\item \emph{Low:} 1-2 dimensions (visualization possible)
\item \emph{Medium:} 3-4 dimensions (moderate curse of dimensionality)
\item \emph{High:} 5+ dimensions (severe curse of dimensionality)
\end{itemize}

\textbf{Budget Dimension:}
\begin{itemize}
\item \emph{Low Budget:} $< 20$ evaluations per dimension
\item \emph{High Budget:} $\geq 20$ evaluations per dimension
\end{itemize}

\subsection{3D Performance Calibration}

For each optimizer $i$ and context $\mathbf{c} = (l, d, b)$, we estimate context-specific performance:

\begin{equation}
\hat{\theta}_{i}(l, d, b) = \frac{1}{N_{i,l,d,b}} \sum_{k \in \mathcal{P}_{l,d,b}} R_{i,k}
\end{equation}

where $R_{i,k}$ is optimizer $i$'s normalized rank on problem $k$, and $\mathcal{P}_{l,d,b}$ is the set of problems matching context $(l, d, b)$.

This produces a 3D performance tensor $\Theta \in \mathbb{R}^{|\text{Optimizers}| \times 3 \times 3 \times 2}$ enabling context-specific recommendations.

\section{Browser-Based Implementation} \label{sec:implementation}

\subsection{Zero-Installation Architecture}

\pkg{HumpDay} prioritizes accessibility through browser-native execution using Pyodide, eliminating installation barriers that limit adoption of optimization tools:

\begin{lstlisting}[language=JavaScript]
// Load Python environment in browser
pyodide = await loadPyodide();
await pyodide.loadPackage(['numpy', 'scipy']);

// Execute 3D calibration analysis
pyodide.runPython(`
    from humpday.calibration import ThurstoneCaliber3D
    caliber = ThurstoneCaliber3D()
    recommendations = caliber.get_recommendations(
        landscape_type='multimodal',
        dimensionality='high',
        budget='low'
    )
`);
\end{lstlisting}

\subsection{Interactive Context Explorer}

The platform provides an interactive interface for exploring the 3D performance landscape:

\textbf{Context Selection:} Users specify their problem characteristics via intuitive controls.

\textbf{Real-Time Recommendations:} The system immediately updates optimizer recommendations based on selected context.

\textbf{Performance Visualization:} 3D plots show how optimizer performance varies across context dimensions.

\textbf{Confidence Intervals:} Recommendations include uncertainty estimates based on available data.

\section{Experimental Validation} \label{sec:results}

\subsection{Experimental Design}

We conducted 432 optimization runs across:
\begin{itemize}
\item \textbf{3 optimizers:} SciPy implementations (L-BFGS-B, Powell, SLSQP)
\item \textbf{12 objectives:} Spanning smooth, multimodal, and rugged landscapes
\item \textbf{3 dimensionalities:} 2D, 4D, 8D problems
\item \textbf{2 budget levels:} 25 and 75 function evaluations
\item \textbf{2 repetitions:} For statistical reliability
\end{itemize}

All runs achieved 100\% success rate, providing complete performance data across all context combinations.

\subsection{3D Performance Tensor Results}

The multi-dimensional analysis reveals clear specialization patterns:

\textbf{L-BFGS-B Specialization:}
\begin{itemize}
\item Smooth landscapes: 0.694-0.806 performance across all dimensions
\item High-dimensional smooth problems: 0.806 (best performance observed)
\item Multimodal landscapes: 0.500-0.736 (context-dependent)
\end{itemize}

\textbf{Powell Method Specialization:}
\begin{itemize}
\item Rugged landscapes with high budget: 0.917 (highest single performance)
\item Multimodal landscapes: 0.611-0.722 performance
\item Smooth landscapes: 0.694 (competitive but not dominant)
\end{itemize}

\textbf{SLSQP Characteristics:}
\begin{itemize}
\item Highest consistency across contexts (standard deviation 0.072)
\item Lower peak performance but reliable across conditions
\item Overall versatility ranking: 1st (combined metric)
\end{itemize}

\subsection{Context-Specific Recommendation Validation}

The 3D calibration produces actionable recommendations:

\begin{center}
\begin{tabular}{lll}
\toprule
Context & Recommended Optimizer & Performance \\
\midrule
Smooth, High-Dim, Low Budget & L-BFGS-B & 0.806 \\
Rugged, Low-Dim, High Budget & Powell & 0.917 \\
Multimodal, Low-Dim, High Budget & Powell & 0.722 \\
Smooth, Medium-Dim, Low Budget & L-BFGS-B & 0.750 \\
\bottomrule
\end{tabular}
\end{center}

These recommendations demonstrate substantial performance differences based on context, with the best choice varying by up to 0.4 performance units.

\subsection{Comparison with Global Rankings}

Traditional global ranking would suggest:
\begin{enumerate}
\item Powell (0.603 overall)
\item L-BFGS-B (0.571 overall)
\item SLSQP (0.492 overall)
\end{enumerate}

However, context-specific analysis reveals this global ranking misleads in many scenarios:
\begin{itemize}
\item L-BFGS-B outperforms Powell by 0.112 on smooth high-dimensional problems
\item Powell outperforms L-BFGS-B by 0.417 on rugged high-budget problems
\item SLSQP provides the most reliable performance across diverse contexts
\end{itemize}

\section{Educational and Practical Impact} \label{sec:impact}

\subsection{Interactive Learning Framework}

The browser-based platform enables interactive exploration of optimization concepts:

\textbf{Algorithm Personalities:} Each optimizer is presented with characteristic ``strengths'' and ``preferences'' based on empirical performance patterns.

\textbf{Problem Landscape Visualization:} Real-time 3D surface plots show how different optimizers navigate various landscape types.

\textbf{Context Sensitivity Demonstrations:} Users can adjust problem parameters and observe how recommendations change, building intuition about context dependency.

\subsection{Practical Decision Support}

For practitioners, the 3D calibration framework transforms optimizer selection:

\textbf{Before:} ``I'll try a few popular methods and see what works''

\textbf{After:} ``Based on my smooth 10-dimensional problem with 50-evaluation budget, L-BFGS-B is recommended with 0.806 expected relative performance''

This evidence-based approach reduces trial-and-error and improves optimization outcomes.

\section{Discussion and Future Work} \label{sec:discussion}

\subsection{Implications for Optimization Practice}

The multi-dimensional Thurstone calibration approach addresses a fundamental gap between optimization research and practice. While the research community often focuses on algorithmic improvements measured on standard benchmarks, practitioners need guidance for their specific problem contexts.

Our results demonstrate that the question ``Which optimizer is best?'' has no universal answer—but the question ``Which optimizer is best for my problem type?'' does have principled, data-driven answers.

\subsection{Scalability and Extensibility}

The current 3D framework (landscape × dimensionality × budget) represents a minimal viable analysis. Natural extensions include:

\textbf{Additional Dimensions:} Noise characteristics, constraint types, objective function smoothness measures

\textbf{Finer Granularity:} More nuanced problem classifications based on landscape analysis

\textbf{Dynamic Calibration:} Online updating of performance estimates as new benchmark data becomes available

\textbf{Confidence-Weighted Recommendations:} Incorporating uncertainty estimates when limited data exists for specific contexts

\subsection{Community Benchmarking Platform}

The browser-based architecture enables crowdsourced benchmark contributions:
\begin{itemize}
\item Researchers can submit new optimizers via standardized API
\item Problem contributors can add new objective functions
\item Performance data aggregates across multiple contributors
\item Reproducible evaluation through deterministic seeding
\end{itemize}

\section{Conclusion} \label{sec:conclusion}

We have presented multi-dimensional Thurstone calibration as a principled approach to context-specific optimizer recommendation. Key contributions include:

\begin{enumerate}
\item \textbf{Theoretical Framework:} Extension of Thurstone models to multi-dimensional optimization performance analysis
\item \textbf{Empirical Validation:} Demonstration that optimizer performance exhibits strong context dependency across landscape type, dimensionality, and budget dimensions
\item \textbf{Practical Tool:} Browser-based platform providing zero-installation access to context-specific recommendations
\item \textbf{Educational Resource:} Interactive framework for understanding optimization algorithm behavior
\end{enumerate}

The transformation from global rankings to context-specific recommendations represents a significant advance in making optimization research accessible and actionable for practitioners. By asking not ``Which optimizer is best?'' but ``Which optimizer is best for problems like mine?'', we enable evidence-based algorithm selection that improves optimization outcomes across diverse application domains.

The open-source platform at \url{https://microprediction.github.io/humpday} demonstrates that sophisticated optimization analysis can be democratized through browser-based implementation, eliminating installation barriers while maintaining statistical rigor.

Future work will extend this framework to additional problem dimensions, incorporate more optimization algorithms, and develop adaptive recommendation systems that learn from user feedback to improve context-specific guidance over time.

\section*{Software Availability}

\pkg{HumpDay} is freely available under the MIT license at \url{https://github.com/microprediction/humpday}. The browser-based demo with 3D Thurstone calibration is accessible at \url{https://microprediction.github.io/humpday} without installation requirements.

Installation via package manager:
\begin{lstlisting}[language=bash]
pip install humpday
# Or with browser-compatible optimizers only:
pip install humpday[browser]
\end{lstlisting}

All experimental data and analysis scripts are included in the repository under \texttt{experiments/}, enabling full reproduction of results.

\section*{Acknowledgments}

The author thanks the optimization community for feedback on early versions of this work and the Pyodide project for enabling Python in browsers. Special thanks to the Thurstone package developers for statistical modeling foundations.

% Bibliography
\bibliography{humpday-jss}

\end{document}